\documentclass[../main.tex]{subfiles}

\graphicspath{{../image/}{../figure/}}

\begin{document}

\chapter{两个守恒律的真非线性方程组}\label{chapter:ch12-两方程真非线性组}

本章将以与第 XI 章中所述真非线性单个守恒律相当的精度，建立两方程真非线性严格双曲守恒律组解的理论。
这将借助 Riemann 不变量坐标系的存在以及由此产生的丰富的熵--熵流对族来实现。本研究中的主要工具将是广义特征线与熵估计。

分析将揭示单个守恒律解与成对守恒律解在结构上的密切相似性。因而，与单个情形一样，跳跃间断一般由激波的碰撞和/或压缩波的聚焦产生，随后被分解为波扇，而波扇由相应 Riemann 问题的解局部逼近。

解沿类空曲线的迹的总变差由初始数据的总变差控制，而稀疏波的铺展影响总变差，这与单个情形相同。

在单个情形中遇到的耗散机制在此同样起作用，并对解的大时间行为具有类似的效应。熵耗散诱导初始数据属于 $L^1(-\infty,\infty)$ 的解以 $O(t^{-1/2})$ 衰减。当初始数据具有紧支集时，两个特征族渐近退耦；特征线铺展，每个族形成单个 $N$-波剖面。最后，与单个情形一样，周期初始数据下特征线的禁锢诱导每个周期总变差以 $O(t^{-1})$ 衰减，并形成锯齿状剖面，每个特征族一个。

\subfile{Chapter-12/section-01.tex}
\subfile{Chapter-12/section-02.tex}
\subfile{Chapter-12/section-03.tex}
\subfile{Chapter-12/section-04.tex}
\subfile{Chapter-12/section-05.tex}
\subfile{Chapter-12/section-06.tex}
\subfile{Chapter-12/section-07.tex}
\subfile{Chapter-12/section-08.tex}
\subfile{Chapter-12/section-09.tex}
\subfile{Chapter-12/section-10.tex}
\subfile{Chapter-12/section-11.tex}

\end{document}
